\section{Input data for operating mode} \label{sec2.1}

\medskip
\textbf{General Remarks}

\medskip

The variables in this block control some of the more general options
in EIRENE such as overall running time, usage of dump files but also
a few parameters depending on
the simulation model (drifts included or not, etc.).

\medskip

\textbf{The Input Block}

\begin{lstlisting}
      TXTRUN
*** 1. Data for operating mode
\end{lstlisting}

first card:  integer flags controlling global options for the entire run.

\begin{lstlisting}
      NPRLL  NMODE  NTCPU  NFILE  NITER0  NITER  NTIME0  NTIME
\end{lstlisting}

next input line is only for EIRENE$_{2002}$ and younger, and optional in these
versions. Defaults, if this card is not included, are given below.

\begin{lstlisting}
      NOPTIM NOPTM1 NGEOM_USR NCOUP_INPUT NSMSTRA NSTORAM
     .  NGSTAL NRTAL NREAC_ADD
\end{lstlisting}

next line (second mandatory line): logical flags for global options [format: 12(5L1,1X)].

\begin{lstlisting}
      NLSCL NLTEST NLANA NLDRFT NLCRR
     .  NLERG NLIDENT NLONE NLMOVIE NLDFST
     .  NLOLDRAN NLCASCAD NLOCTREE NLWRMSH
\end{lstlisting}

next input lines are only for EIRENE$_{2005}$ and younger, and
optional in these versions. There can be any number of these
cards starting with character string \lstinline|CFILE|  in the input file.

\begin{lstlisting}
      "CFILE" DBHANDLE DBFNAME
\end{lstlisting}

\medskip

\textbf{Meaning of the input flags for operating mode}

\medskip

\begin{description}
\item[TXTRUN]
At least one line of text to identify the run
on printout and plot files.

Each text card must start with *. The first of these text lines
will appear on all plots, the full text will be echoed at the
beginning of the printout file
\item[NPRLL]
(Formerly at this place: NMACH, flag for machine precision (round off errors), now: out.)

The new (2018) flag NPRLL controls the various options for (MPI) parallelization modes, i.e.\ the strategies
to assign processes  to strata. This is defined in PEDIST.f (initialization phase)
by filling the matrix PROCFORSTRA(strata,processors).
\begin{description}
\item[$= 1$] old default: attempting proportional allocation and, simultaneously, load balancing (see under: stratified sampling,
the remarks towards the end of \cref{sec1.2.1b}).
\item[$= 2$] not in use
\item[$= 3$] not in use
\item[$= 0$] new default: ``embarrassingly parallel mode'': all processes do the full MC job (all strata).

Most transparent, and simplest method, load balancing guaranteed. But perhaps not optimal wrt.\ communication
\item[$= -1$] user defined distribution of strata to processes (subr.\ PEDIST, calls: PEDIST\_USR.f)
\end{description}



\item[NMODE] Operating mode
\begin{description}
\item[$= 0$] EIRENE run as stand-alone-code. Code coupling segment
\lstinline|couple_Dummy.f| may be used. 
Input block 14 has a fixed format, see \cref{sec2.14}.
\item[$\ne 0$]
EIRENE calls the code coupling subroutine INFCOP in the code coupling segment
%\begin{lstlisting}
\lstinline|couple_Name.f|
%\end{lstlisting}
where 'Name' is a character string identifying the
particular external code, to which EIRENE is to be coupled
(e.g.: Name = B2, Name = DIVIMP, Name = U-file, Name = FIDAP, Name = EMC3,
Name = OSM, etc.).

Hence: the routine INFCOP
is called by EIRENE for communication with external data sources,
e.g., external data structures for iterative mode by coupling to other codes.
The calling program for the first 3 entries of INFCOP
is subroutine INPUT, and the call to subroutine INFCOP is
after reading 13 blocks from the formatted input file (READ (IUNIN, \dots)).
A 14$^{th}$
block of the input file is read from subroutine INFCOP with the format
(if any) specified there.
At entry IF0COP geometrical
data are provided (overruling corresponding data in input block 2).

At entry IF1COP plasma profiles
(more generally: background medium data) are defined (overruling the
background data specified in block 5).

Any other data (e.g., surface- or volume source distributions overruling
the data in input block 7)
are expected from entry IF2COP.

\item[$> 0$]
If NMODE $> 0$, subroutine INFCOP is called once again after each
completed stratum (at the entry IF3COP(ISTRAA,ISTRAE)) to return
data to another code (post-processing). The call to IF3COP is from
subroutine MCARLO, at the end of the DO 1000 -loop over the strata. One
final call to the interfacing routine (entry IF4COP) is
implemented after the calls to the EIRENE printout- and plot
routines, e.g.\ to perform global balances etc.\ after summation of
the contributions from the individual strata.

\end{description}
\item[NTCPU]
Maximum number of CPU seconds allowed for this run. NTCPU must be
less than or equal to the time parameter in the job-card (if any).

If more than one iteration is carried out (NITER, see below, this
section) or more
than one time cycle (NTIME, see below, this section), then each
iteration or cycle can take up to NTCPU seconds.

In 2008 the definition of NTCPU was slightly changed:
the initialisation overhead is not any longer included in NTCPU. I.e.,
this variable NTCPU now is close to the true Monte Carlo sampling time.
The total CPU-time, including initialisation overhead as well as post-processing
is printed at the end of an EIRENE run (see: ``total CPU-time of this run"
in printout file).
\item[NFILE]
Flag for the use of dump files FT10, FT11, FT12, FT13, FT14 and FT15.

= 0: Neither reading nor saving of data is done.

= JKLMN   (NFILE a 5 digit integer
\begin{description}
\item [N = 1 ] EIRENE writes output data from this run into
files
FT10 and FT11 to save them for plot or printout options in this or in later
'read only' runs with the same input file.
\item [N = 6 ] same as NFILE-N=1, but only the data for the sum over
all strata are written (sufficient, e.g., for \acs{BGK}-iterations, or for post-processing,
(subroutine DIAGNO, see input block 12)
as long as this post-processing only is to be done on total results (not on
contributions from individual strata).
\item [N = 2 ]
EIRENE reads output data from an earlier run from files FT10 and
FT11. No new particle histories are computed, only the requested
output is printed and plots are produced. Iteration cycles or time
cycles (if any) are abandoned.
\item [N = 7 ] same as NFILE-N=2, but only the data for the sum over strata
are read. See N=6 option described above.
\item [M = 1 ]
EIRENE writes geometry data into file FT12. These data are the
output from subroutines GRID and VOLUME in
the initialization
phase.
\item [M = 2 ]
EIRENE reads geometry data from file FT12. The geometry subroutines
GRID and VOLUME are not called. This option should be used if the
geometry has not changed as compared to an earlier run. It reduces the
CPU costs for the overhead.
\item [L = 1 ]
EIRENE writes plasma data and atomic and molecular
data (``A\&M data") into file FT13.
These data are output from subroutines
PLASMA, XSECTA, XSECTM, XSECTI, XSECTP
in the initialization phase. At the end of a run,
some background medium data may have been modified in subroutine MODUSR, (iterative mode) see NITER
flag below. In this case FT13 may be re-written to prepare for next iterations.

Also the primary source parameters (input block 7) are
saved, for later iterations or time-steps.
\item [L = 2 ]
EIRENE reads plasma data, A\&M data from file FT13.
It also reads the entire input block 7 from FT13, and overwrites the input read from this block 7
by those data.
Routines PLASMA,
XSECTA, XSECTM, XSECTI, XSECTP  are not called.
This option should be used if neither the plasma background
nor the selection of atomic processes to be used,
nor the primary source model (block 7)
has changed as compared to an earlier run. It then reduces the
CPU costs for the overhead.
\item [L = 3 ]
Acts as if both NFILE-L=1 and NFILE-L=2. I.e.\ reading background data from file, and writing new
background data onto file at the end of the run, the time-step or the iteration.
For continuation of iterative calculations, for example.
\item [L = 4 ]
Same as NFILE-L=3, except: primary source data (input block 7) are not
read from fort.13, but are taken as specified in input block 7.
This allows to modify these primary sources parameters during iterations or
time-steps not only via module MODUSR, but directly via input file.
\item [L = 6,7,8,9 ]
Same as L=1,2,3,4, respectively, but XDR file format is used.
\item [K = 1 ]
EIRENE saves some data for optimizing non-analog sampling and
stratified source sampling on file FT14.
\item [K = 2 ]
EIRENE reads from file FT14 and tries to optimize operation for the next run, time-step or iteration. Currently only
allocation of CPU time in stratified source sampling is optimized. More details: see \cref{NFILEK} below.
\item [K = 3 ]
Acts as if both NFILE-K=1 and NFILE-K=2.
\item [J = 1 ]
Only for time-dependent option (see block 13). EIRENE writes
``census data" at the end of last time-step onto file FT15.
\item [J = 2 ]
EIRENE reads ``census data" from file FT15, and uses it as initial
distribution for the coming next time-step.
\item [J = 3 ]
Acts as if both NFILE-J=1 and NFILE-J=2.
\end{description}
\item[NITER0]
Initial iteration number: (Default: NITER0=1)
Irrelevant parameter. Only needed
for book keeping and printout. EIRENE labels the iterations from NITER0 to NITER.
\item[NITER]
Number of iterations, if EIRENE runs in ``iterative mode".
\begin{description}
\item[$> 0$] EIRENE calls user supplied subroutine MODUSR
after completing the run;
some model parameters may be modified here for the next iteration step,
and some results from the previous step may be saved on a file.
\item[$> 1$]
EIRENE recalls itself but does not read from the
formatted input file again. This recalling is repeated NITER times
(including the first iteration). The CPU time NTCPU
is used for each iteration. Hence the true CPU time then is
$ NTCPU \cdot NITER $.
\end{description}
\item[NTIME0]
Initial time-cycle number: (default: NTIME0=1). Irrelevant. Only needed
for book keeping and printout. EIRENE labels the time steps from NTIME0 to NTIME.
\item[NTIME]
Total number of iterations (``cycles") in time carried out in one single run.
The total time
per cycle is defined as NTMSTP $\cdot$ DTIMV (see below, input block 13).

After each time-cycle the subroutine TMSTEP is called.
In this routine the ``census arrays", which store the test particle
population at time $t_i$:

$t_{i} = t_{i-1} + \Delta t = t_0 + i \Delta t = TIME0 + ITIME
\cdot [NTMSTP \cdot DTIMV]$

are filled and
prepared for the next time-cycle.
The census arrays from the previous time cycle (if any), i.e., at $t =
t_{i-1}$, are overwritten here.

After the last time-cycle, the census arrays are written on file
fort.15, in order to permit continuation in time in a next run.

The background
conditions (and source distributions or any other input parameters)
for the next time cycle can be modified in subroutine
TMSUSR($t_i$), which is called from subroutine TMSTEP.

If the population on the census array is not empty or known from a
previous run (NFILE-J flag, see above), then this census
population defines one additional stratum for the current cycle.
I.e., the census population then determines the initial condition
for the distribution function
$f(\underline{r},\underline{v},i,t=t_0)$. The source strength FLUX
is computed from that initial condition.
\item[NOPTIM]
Default: NRAD, = total number of grid cells

(verify this default in case of older versions: there it may have been NOPTIM=1)

NOPTIM is the first dimension of the arrays IGJUM3(ICELL,ISURF), which may be
used for optimizing code performance by reducing unnecessary geometrical calculations.
See `CH3-cards' in input block 8 (\cref{sec2.8}).
\item[NOPTM1]
Default: 1
\item[NGEOM\_USR]
for value =1: user defined (external) geometry package (LEVGEO=10),
then no grid storage is provided in EIRENE.
Default: 0
\item[NCOUP\_INPUT]
=1: Storage for data transfer via coupling routines, =0: no
such storage.
Default: 1
\item[NSMSTRA]
``Sum over strata" disabled
for value 0, enabled for value 1. Default: 1
\item[NSTORAM]
Storage vs.\ speed in atomic data evaluation. Maximum
storage, fastest computation: =9. Minimum storage, maximum work
(slowest option) =0. Default: 9
\item[NGSTAL]
Storage for spatial distribution of surface tallies on non-default
standard surfaces for value =1. For value =0: only total (spatially
integrated) surfaces fluxes.
Default: 0
\item[NRTAL]
Condensing mesh cells into fewer larger
cells, for output volume tallies. The underlying fine mesh has NRAD cells (see input block 2).
The coarser mesh, obtained from condensing cells into one larger cell, has NRTAL cells.
Default: 0: Then internally: NRTAL=NRAD, and no condensation is carried out.

NCLTAL(IRAD) = IRTAL: grid cell IRAD is condensed into the larger
cell IRTAL. I.e.: output is average over a larger cell IRTAL, which is comprised of
all cells IRAD such that NCLTAL(IRAD) = IRTAL,  IRAD=1, NRAD

The input data in input block 5 (background medium) are always given on
the fine mesh (size: NRAD)

For output tallies (cell averaging) several cells IRAD1, IRAD2, \dots\ 
can be condensed into one larger cell IRTAL. Cell volumes, scoring,
statistics are automatically done on the coarser grid (size: NRTAL).


The index array NCLTAL(NRAD) can be defined in the problem
specific ``user" routines (see \cref{sec3}), e.g.: INIUSR, GEOUSR, etc..

Default: NCLTAL(IRAD)=IRAD for IRAD=1,NRAD
\item[NREAC\_ADD]
Storage for additional reaction decks read onto EIRENE
arrays in USR-rou\-tines, post-processing, etc. Default: 0
\end{description}
Next input card, global logical switches. The switches 10--15 in this card
may have version specific meanings and are for internal code options only.
\begin{description}
\item[NLSCL]
Some volume averaged tallies are re-scaled in
order to exactly preserve the total number of particles, which
otherwise would be the case only up to statistical precision (due
to the use of track-length estimators). EIRENE computes three
factors FATM, FMOL and FION such that particle balances for atoms,
molecules and test ions, respectively, are accurately observed, if
NLSCL = TRUE.
\item[NLTEST]
Tests for consistency
between cell numbers and geometrical data along the particle
tracks are carried out at each point of collision. If
inconsistencies are detected, the history is stopped and an error
message is printed. The contribution of these particles to the
particle- and energy balances is stored in the bins ``PTRASH" and
``ETRASH" respectively.
\item[NLANA] De-activates
(NLANA=.TRUE.) all non-analog sampling distributions, such as
biased source sampling, splitting, etc. Select NLANA=.TRUE., if
particle trajectory plots are used to get an intuitive picture of
what is going on physically.
\item[NLDRFT] Drift
component in the bulk ions velocity distribution is included, i.e.\ the assumed underlying distribution in velocity space
is a drifting Maxwellian for volumetric background tallies of bulk particles.
(see input block 5, input tallies VXIN, VYIN, VZIN)). Otherwise (if NLDRFT = FALSE) an
isotropic Maxwellian distribution is
assumed for the background particles and the input for VXIN, \dots, VZIN is ignored.

This flag also affects the output tallies for energy exchange, momentum exchange between test particles and background particles,
as well as sampling from linear collision kernels.
\item[NLCRR]
Correlated sampling is used. See discussion at end of \cref{sec1.7}.
\item[NLERG]
The case is
automatically reduced to a case for estimating cell volumes by
utilizing an ``ergodic property". More details: see \cref{NLERG} below.
\item[NLIDENT]
In multi-processor calculation mode NLIDENT forces all processors working on the
same stratum to use the same sequence of random numbers and hence to
carry out exactly identical work.
This flag can be used to test the parallelized code version.
\item[NLONE] The case
is automatically reduced to a single species and ``one speed
transport" problem
(to simplify setting up cases for comparison with analytic results). (not ready, don't use)
\item[NLMOVIE]
reset a number of model parameters to enable a series of geometry
plots for making particle trajectory movies. This option only works in connection with time dependent mode, see NTIME flag
described above, and input block 13 (\cref{sec2.13}). More details: see \cref{NLMOVIE} below.
\end{description}
The next 5 logical flags are for internal use only, in proprietary parts of the code.
They should not be activated without checking the underlying code.
Note in particular:  NLOLDRAN, see below.
These variables may disappear and re-appear with other names/features,
no code backward compatibility. They should all be set to FALSE, except probably NLOLDRAN,
with may be code version dependent.
\begin{description}
\item[NLDFST]
for internal use. Testing time step options in Fokker Planck solver (folion.f)
\item[NLOLDRAN]
Switch between random number generators. Opposite meaning in different code versions.
One should set this switch such that generator ``H1RN" is used, and NOT the congruential generator.
See function ``ranf.f". The ``junk" congruential generator has been kept, so far,
to permit full backward compatibility. But the code now (from revisions April 17 on) heavily complains if it is used.
\item[NLCASCAD]
for internal use. Unfinished option for enforcing fully analog collision cascades,
rather than weighting  (collide.f)
\item[NLOCTREE]
for internal use. Testing CPU optimization for timea.F routines (``additional surfaces")
\item[NLWRMSH]
write out additional surface information for internal (FZJ) triangular grid generators (input.f)
\end{description}
Optional input cards, for pathways to external databases.
Their presence is identified by the starting character string CFILE of these cards.
\begin{description}
\item[DBHANDLE]
Name of the database which is specified in this ``CFILE" card.
The format of a data file is identified in EIRENE by its name.
Possible database names currently recognized are:

AMJUEL, HYDHEL, METHAN, H2VIBR, HYDRTC, ADAS,  SPECTR,
PHOTON, POLARI,
\medskip

SPUTER, TRIM,  TRMSGL,

with volumetric collision databases in the first block, and
surface interaction databases in the second.
\item[DBFNAME]
Absolute or relative path and name of file to be used for this
database.
\end{description}
Example:

\lstinline|CFILE AMJUEL /home/path/AMdata/amjuel.tex|

\lstinline|CFILE TRIM /home/path/Surfacedata/TRIM/trim.dat|

Default:
If the specification of an external database name and path,
which is later (e.g.\ in block 4 (\cref{sec2.4}) or block 6 (\cref{sec2.6})) referred to in
a particular run, is not given here, then
these data files are expected, usually with the same name,
in the local directory from which the calculation
takes place. Due to historical reasons and for backward compatibility for some
data files also different default names are recognized.

Default file names to be used in this local directory currently are

AMJUEL,

HYDHEL,

METHAN or ``METHANE", for database name METHAN

H2VIBR,

HYDRTC,

SPECTR,

PHOTON,

POLARI,

SPUTER,

TRIM or ``fort.21",  for database name TRIM

For ADAS (AMdata) and TRMSGL (Surfacedata) there is an exception from this scheme.
In these cases no default file names are available.
The notation ``ADAS" and ``TRMSGL" here are only synonymous for any atomic or surface data file
which is stored in the same format as the so called \ac{ADAS} adf11 files and the TRIM
conditional quantile data tables, independent on
whether these tables originate from \ac{ADAS} or TRIM.

And input variable DBFNAME
specifies only the path to the directory
which is the root of the data tabulated in \ac{ADAS} (tabular) format and the TRIM format
data tree, respectively.

Example:

\lstinline|CFILE ADAS /home/path/AMdata/Adas_Eirene_2010/adf11/|

\lstinline|CFILE TRMSGL /home/path/Surfacedata/TRIM/|

The particular data file from these trees (e.g.: \lstinline|SCD96| in case of \ac{ADAS}, or \lstinline|A_on_B|
in case of TRMSGL), the
chemical element  and the charge stage of a
particular ion $Z = 1, \dots, Z+1$ (e.g.: C, 4), are detailed later in
input block 4 (for volumetric collision processes based on \ac{ADAS} adf11 format), and input block 6 (TRMSGL, for the surface reflection database based on TRIM)
see below, \cref{sec2.4,sec2.6},
respectively.

\subsection{Automated stratification optimization: proportional allocation}\label{NFILEK}
The stratified sampling in EIRENE can improve the performance (measured in terms of statistical error per CPU-time),
if proportional allocation of sample size $N_k$ to strata source strength $S_k$ can be achieved. See \cref{sec1.2.1b}
for more information on stratified source sampling.
The input flag ALLOC in input block 7 (\cref{sec2.7}) controls this allocation.

Monte Carlo codes such as EIRENE are, in principle,  trivially parallelizable by just distributing the Monte Carlo
trajectory calculation over the available compute nodes.
If parallelization is activated in combination with stratified source sampling, however, and strata are assigned to compute nodes,
then load balancing may be hard to achieve together with proportional allocation of weights to strata. This is because trajectories from
some strata may be very short (e.g.\ transport into purely absorbing media, for example of impurity particles until ionisation) and very fast to compute,
while from other strata these trajectories may be very long and time consuming to compute, e.g.\ in near diffusive regimes with many collisions.
Often volume recombination strata in cold plasma regions fall into this category.

Load balancing and proportional (or otherwise optimal) allocation to strata then requires estimates of run-time per history for each stratum
prior to the Monte Carlo simulation itself.

This information is stored in output stream FORT 14, and can be read back in into the next EIRENE run and used then for
assignment of compute nodes to strata and vice versa. This is controlled by the NFILE-K flag described above.


\subsection{The NLERG option for cell volumes}\label{NLERG}
All non-transparent surfaces are made perfectly reflecting
(mono-energetic, cosine distribution). All volumetric collision
processes are de-activated. Except, possibly, for the very first
flight, there is only either one atomic species (IATM=1) or one
molecule species (IMOL=1). Unless otherwise specified in input
block 13, a single time-step, with time limit DTIMV=\num{1.0} (\si{\second}) is set
and up to NPRNLI=10 test flights are launched. Under such
conditions, the particle density should be constant throughout the grid. Statistically
significant deviations from that constancy indicate wrong cell
volumes as compared to the volumes seen by the test particles
(e.g.\ due to additional surfaces intersecting the grids, or
complicated shapes of some additional cells (see input blocks 2
and 3). In a second run the volumes of these cells can then be set explicitly, e.g.
in input block 8.
\subsection{The NLMOVIE option for making movies of trajectories}\label{NLMOVIE}
documentation to be written
